\chapter[选择研究参与者：界定、抽样与招募]{选择研究参与者：\\界定、抽样与招募}
\chapterauthors{Warren S. Browner、Thomas B. Newman、Mark J. Pletcher}
% Chinese reconstructions of source PDF figures 3.1--3.4.
\newcommand{\chthreeoverview}{%
\begin{tikzpicture}[x=1cm,y=1cm,font=\small,>=Stealth,
 panel/.style={draw=dcrblue!65,line width=.7pt,fill=white,align=center,text width=3cm},
 base/.style={draw=dcrred,fill=white,line width=.8pt,minimum width=3.15cm,minimum height=.9cm,align=center}]
\foreach \x/\t/\a/\b in {0/研究问题/目标人群/目标现象,5/研究计划/预定样本/预定变量,10/实际研究/实际样本/实际测量}{
\node[panel,minimum height=2.5cm] (c\x) at (\x,0) {};
\node[fill=dcrblue!65,text=white,align=center,minimum width=3.15cm,minimum height=.7cm] at (\x,1.05) {\bfseries\t};
\node at (\x,.25) {\a};\node at (\x,-.65) {\b};}
\draw[->,line width=3pt,dcrblue!65] (1.6,.25)--(3.4,.25) node[midway,above,font=\small] {设计};
\draw[->,line width=3pt,dcrblue!65] (6.6,.25)--(8.4,.25) node[midway,above,font=\small] {实施};
\draw[red,line width=1.2pt] (-1.75,-.15) rectangle (11.75,.65);
\node[text=dcrgold,align=center] (choice) at (5,-2.2) {选择与误差\\影响……};
\node[text=gray,align=center] (ev) at (2,-3) {外部\\效度};
\node[text=gray,align=center] (iv) at (8,-3) {内部\\效度};
\node[base] (understand) at (0,-5) {现有认识};
\node[base] (conclusion) at (5,-5) {研究结论};
\node[base] (result) at (10,-5) {研究结果};
\draw[->,gray,thick] (choice)--(ev);\draw[->,gray,thick] (choice)--(iv);
\draw[->,gray,thick] (ev)--(2.5,.1);\draw[->,gray,thick] (ev)--(understand);
\draw[->,gray,thick] (iv)--(7.5,.1);\draw[->,gray,thick] (iv)--(11.85,-3);
\draw[->,gray,thick] (iv)--(conclusion);
\draw[->,line width=3pt,dcrred!75] (result.west)--(conclusion.east) node[midway,above,font=\small] {解释};
\draw[->,line width=3pt,dcrred!90] (conclusion.west)--(understand.east) node[midway,above,font=\small] {推断};
\draw[->,line width=2pt,dcrblue!70] (c10.east)--(12.3,0)--(12.3,-5)--(result.east) node[pos=.43,right,font=\small,text=black] {分析};
\draw[->,line width=2pt,dcrred!85] (understand.west)--(-2.3,-5)--(-2.3,0)--(c0.west) node[pos=.43,left,font=\small,text=black] {思考};
\end{tikzpicture}}

\newcommand{\chthreegeneralize}{%
\begin{tikzpicture}[font=\footnotesize,>=Stealth,
 p/.style={draw=dcrblue!70,fill=white,align=left,text width=2.65cm,inner sep=5pt},
 b/.style={draw=dcrred,align=center,minimum width=2.7cm,minimum height=.8cm}]
\node[p] (target) at (0,0) {\textbf{目标人群}\\\textbf{推广较有把握：}\\研究期间美国所有郊区白人具有同样的关联。\\[4pt]\textbf{推广把握较小：}\\同样的关联存在于：\\其他成人（如黑人）；\\儿童；其他国家的人；\\2050 年的人。};
\node[p] (accessible) at (4,.25) {\textbf{可接触人群}\\弗雷明汉所有成人中存在同样的关联。};
\node[p] (intended) at (8,.25) {\textbf{预定样本}\\按设计选出的弗雷明汉成人样本中存在同样的关联。};
\node[p] (actual) at (12,.25) {\textbf{实际样本}\\在研究参与者中观察到高血压与冠心病之间存在关联。};
\draw[->,dcrblue!70,line width=2pt] (actual)--(intended);
\draw[->,dcrblue!70,line width=2pt] (intended)--(accessible);
\draw[->,dcrblue!70,line width=2pt] (accessible)--(target);
\node[text=dcrgold,align=center] (choice) at (6,-3) {选择与误差\\影响……};
\node[text=gray,align=center] (ev) at (2.5,-3.4) {外部效度};
\node[text=gray,align=center] (iv) at (9.5,-3.4) {内部效度};
\node[b] (u) at (0,-5.5) {现有认识};\node[b] (c) at (6,-5.5) {研究结论};\node[b] (r) at (12,-5.5) {研究结果};
\draw[->,gray,thick] (choice)--(ev);\draw[->,gray,thick] (choice)--(iv);
\draw[->,gray,thick] (ev)--(2,-.1);\draw[->,gray,thick] (ev)--(u);
\draw[->,gray,thick] (iv)--(6,-.1);\draw[->,gray,thick] (iv)--(10,-.1);\draw[->,gray,thick] (iv)--(c);\draw[->,gray,thick] (iv)--(13.5,-3.4);
\draw[->,dcrred!70,line width=2pt] (r)--(c) node[midway,above,text=black] {解释};
\draw[->,dcrred!90,line width=2pt] (c)--(u) node[midway,above,text=black] {推断};
\draw[->,dcrblue!70,line width=2pt] (actual.east)--(14,0)--(14,-5.5)--(r.east) node[pos=.45,right,text=black] {分析};
\draw[->,dcrred!85,line width=2pt] (u.west)--(-2,-5.5)--(-2,0)--(target.west) node[pos=.45,left,text=black] {思考};
\end{tikzpicture}}

\newcommand{\chthreesteps}{%
\begin{tikzpicture}[>=Stealth,font=\small,
 chthreeStep/.style={draw=dcrblue!70,fill=dcrlight,text width=3.5cm,minimum height=3cm,align=left,inner sep=7pt}]
\node[chthreeStep] (a) at (0,0) {\textbf{步骤 1：目标人群}\\[7pt]明确适合研究问题的人口学及临床特征。};
\node[chthreeStep] (b) at (5.4,0) {\textbf{步骤 2：可接触人群}\\[7pt]明确时间及“地理”特征，使其既代表目标人群，又可供研究。};
\node[chthreeStep] (c) at (10.8,0) {\textbf{步骤 3：预定样本}\\[7pt]明确抽样方法，选出能够代表可接触人群、愿意且能够充分参与的人。};
\draw[->,line width=3pt,dcrblue!70] (a)--(b);\draw[->,line width=3pt,dcrblue!70] (b)--(c);
\end{tikzpicture}}

\newcommand{\chthreesamples}{%
\begin{tikzpicture}[>=Stealth,font=\scriptsize,
 r/.style={draw=dcrred,align=center,fill=white,text width=2cm,inner sep=3pt}]
\foreach \x in {0,5,10}{\draw[dcrred,dashed,thick] (\x,0) circle (2.2);}
\node[align=center,text=dcrred] (a) at (0,3.2) {接受白内障手术患者的\\可接触人群（$n=2500$）};
\node[align=center,text=dcrred] (b) at (5,3.2) {有子痫前期孕妇的\\可接触人群};
\node[align=center,text=dcrred] (c) at (10,3.2) {收治心肌梗死（MI）患者的\\医院可及总体};
\foreach \x/\n in {0/a,5/b,10/c}{\draw[->] (\n.south) to[bend right] (\x,2.2);}
\node[r] at (0,0) {随机抽取\\100 名患者};
\foreach \i/\y/\t in {1/1.6/拉丁裔,2/.8/黑人,3/0/白人,4/-.8/亚裔,5/-1.6/其他}{
\node[r] (s\i) at (5,\y) {随机抽取\\50 名女性};
\node[anchor=east] (l\i) at (3.2,\y) {\t};\draw[->] (l\i)--(s\i);
}
\foreach \y in {-1.2,-.4,.4,1.2}{\draw[dcrblue] (3.15,\y)--(6.85,\y);}
\draw (9.1,1) circle (.9);\draw (10.95,-.1) circle (.9);\draw (9.1,-1.05) circle (.9);
\node at (9.1,1.72) {医院 A};\node[r,text width=1.5cm] at (9.1,.9) {随机抽取\\75 名 MI 患者};
\node[align=center] at (10.95,-.1) {医院 B\\（不在样本中）};
\node at (9.1,-.25) {医院 C};\node[r,text width=1.5cm] at (9.1,-1.2) {随机抽取\\75 名 MI 患者};
\node[text=dcrblue,font=\small\bfseries] at (0,-2.7) {简单随机样本};
\node[text=dcrblue,font=\small\bfseries] at (5,-2.7) {分层随机样本};
\node[text=dcrblue,font=\small\bfseries] at (10,-2.7) {整群样本};
\end{tikzpicture}}

\sourcepage{58}
研究计划必须明确一个样本：研究它所花费的时间和经费可以接受，人数足以控制随机误差，代表性也足以支持有效解释，并据此推断研究结果在多大程度上可以推广到目标人群。不过，可推广性很少是简单的“是或否”问题，而是一项复杂的定性判断，取决于研究者选择研究谁，以及怎样抽样。

第 5、6 章将讨论如何确定适当的参与者人数。本章介绍如何界定及抽取既有代表性又能实际招募的参与者（图 3.1），并讨论招募他们加入研究的策略。
\sourcepage{59}
\par
\begin{figure}[htbp]\centering
\resizebox{\linewidth}{!}{\chthreeoverview}
\caption*{\textbf{图 3.1}\quad 本章关注红框内的事项：选择一个能够代表研究问题所涉及目标人群的参与者样本。}
\end{figure}\par
并非所有临床研究都以个人为对象；有些研究涉及家庭、工作小组、医院病区、城镇或其他观察单位。不过，本章主要使用以人为参与者的例子。

\section{基本术语与概念}
\subsection{总体与样本}
\term{总体}（population）是具有一组指定特征的所有人的集合。日常用法通常以地理特征界定总体，例如内华达州克拉克县的人口。在研究中，用临床和人口学特征界定\term{目标人群}，即研究以外、希望将结果推广到的较大人群，例如患有哮喘的青少年。\term{可接触人群}则是以地域和时间界定、实际能够研究的目标人群子集，例如某一年在研究者所在社区接受诊疗的哮喘青少年。
\sourcepage{60}
（临床上有时把“患者人群”作为可接触人群的粗略同义词。）

\term{样本}是可接触人群中明确界定的子集。\term{预定样本}（intended sample）是研究者希望纳入的部分，例如在同意开放电子病历的三家医疗中心就诊的哮喘青少年；\term{实际样本}（actual sample）则由真正加入研究的人组成，例如同意入组的人。

\subsection{推广研究发现}
经典的\term{弗雷明汉研究}（Framingham Study）（1）是较早尝试通过样本中的观察，对目标人群作出推断的研究（图 3.2）。
\par
\begin{figure}[htbp]\centering\resizebox{\linewidth}{!}{\chthreegeneralize}
\caption*{\textbf{图 3.2}\quad 研究结束后，把参与者中的结果推广到目标人群，需要作出解释和推断。这个过程从右向左进行。}
\end{figure}\par
\sourcepage{61}
其抽样设计要求先找出马萨诸塞州弗雷明汉所有至少有一名 30～59 岁成员的家庭，按地址列出，再邀请每三个家庭中前两个家庭的适龄人员参加。（这种设计不像逐个随机选择参与者那样不易被人为干预，详见本章后文。）这一方法有几个值得注意之处。最初选中的居民有三分之一拒绝参加，研究者转而接受其他符合年龄条件的自愿者（1）。自愿者通常比拒绝参加者健康，因此实际样本的特征无疑与预定样本不同。更重要的是，弗雷明汉主要由中产阶级白人组成，不能代表全美国，甚至不能代表马萨诸塞州。

不过，任何样本都有一些误差；问题在于这些误差怎样影响根据研究结果所作的解释与推断。弗雷明汉研究的抽样误差似乎并不足以否定以下结论：研究中观察到的风险关系，例如高血压是冠心病（coronary heart disease，CHD）的危险因素，可以推广到弗雷明汉所有居民。下一个问题是，从弗雷明汉这一可接触人群，推广到其他地方的目标人群是否有效。这种推断更具主观性。选中该镇，部分因为它似乎是其他地方类似社区的典型，也因为对研究者方便；研究者所在的大学位于 25 英里之外的波士顿。将观察到的风险关系推广到其他人群，要依赖一组假设，其中许多无法验证。后来研究发现，弗雷明汉居民与许多其他人群，包括内城区黑人，高血压和冠心病的关联相似。但后者高血压更常见。这说明，描述许多特征分布的研究，往往不如研究关联的结果那样容易推广到其他人群。

\sourcepage{62}
\subsection{设计参与者获取方案的步骤}
图 3.2 从右向左展示推断，体现的是解释已完成研究发现时的顺序。规划研究时应倒转这一顺序，从左侧开始（图 3.3）。研究者首先明确最适合研究问题的目标人群临床及人口学特征，再用地域和时间标准界定兼具代表性与可行性的研究样本。
\par
\begin{figure}[htbp]\centering\resizebox{\linewidth}{!}{\chthreesteps}
\caption*{\textbf{图 3.3}\quad 设计研究时，研究者按从左向右的顺序拟订选择参与者的方案。“地理”一词使用得较宽泛，例如也可指某医疗保险计划的成员，甚至某社交媒体平台的用户。}
\end{figure}\par

\section{选择标准}
假设研究者想比较低剂量睾酮补充剂与安慰剂，对提高绝经后女性性欲的疗效。她应首先明确准备把结果应用于什么样的目标人群，然后使样本参与者尽可能接近这一理想人群（表 3.1）。这些特征统称\term{选择标准}或\term{入组标准}（selection or entry criteria）。
\begin{bookbox}{表 3.1\quad 低剂量睾酮与安慰剂改善绝经后性欲的临床试验：选择标准的设计}
\begin{tabularx}{\linewidth}{YY}
\toprule 设计要点 & 例子\\\midrule
\multicolumn{2}{l}{\textbf{纳入标准：具体说明研究谁}}\\
与研究问题直接相关的临床特征 & 担心性欲减退的女性；有性伴侣\\
人口学范围 & 50～59 岁\\
地域（行政）范围、时间范围 & 在指定年份 1 月 1 日至 12 月 31 日期间，到研究者所在医院门诊就诊的患者\\\midrule
\multicolumn{2}{l}{\textbf{排除标准：谨慎设置；不研究以下子集的原因}}\\
很可能失访 & 物质滥用；计划搬到州外\\
无法提供良好数据 & 定向障碍或认知受损；存在语言障碍\textsuperscript{a}\\
发生潜在不良作用的风险较高 & 有心肌梗死或卒中史\\\bottomrule
\end{tabularx}
\smallskip\footnotesize\textsuperscript{a}如果存在语言障碍的亚组人数较多，而且与研究问题相关，可以采集非语言数据，或使用双语工作人员及问卷，以免排除他们。
\end{bookbox}
\sourcepage{63}
\subsection{建立选择标准}
\term{纳入标准}规定与研究问题及设计有关的目标人群主要特征。本例最相关的人群是担心性欲减退的女性；此外，有性伴侣很可能是衡量睾酮补充效果的关键。年龄也往往十分重要。一位研究者可能认为 50 多岁女性的获益与伤害之比最好，因此关注这一年龄层；另一位则可能作出不同选择，关注更年长的女性。
\sourcepage{64}
许多队列研究招募发生研究结局风险高于平均水平的人，以减少所需样本量（见第 5 章）。例如，研究骨质疏松性骨折危险因素时，可以纳入 65 岁及以上女性，因为她们的骨折风险高于年轻女性或男性；骨质疏松性骨折研究（Study of Osteoporotic Fractures）正是采用这一策略（2）。同样，降低结肠癌风险的新疗法试验，可以纳入有结肠息肉病史的参与者。

纳入标准常需要在科学与实际目标之间取舍。研究者所在医院的患者通常容易获得，但当地转诊模式的特殊性，可能妨碍把结果推广到其他人群。确定标准时，应作出合理、能在整个研究中一致应用，而且能向他人清楚说明的决定，让读者能够判断已发表结论适用于谁。有时，还应考虑怎样用纳入标准纠正过去使特定群体代表不足的偏倚或决定。例如，骨质疏松性骨折研究最初只纳入风险最高的白人女性，而后续队列也招募了黑人女性及男性（3、4）。

\term{排除标准}是在宽泛纳入标准基础上规定的有限例外，应谨慎使用。通常只应排除那些可能影响知情同意能力、数据质量、随机治疗可接受性或成功随访的特征（表 3.1），例如英语交流困难、物质滥用等严重心理问题，以及严重疾病。与观察性研究不同，临床试验更可能因为担心某些参与者的干预安全性而设置排除条件，例如妊娠或可能妊娠的女性。

当然，某些纳入标准很容易改写为排除标准，反之亦然。例如，“70 岁以上且不吸烟”，可写作“不纳入 69 岁及以下者或吸烟者”。
\sourcepage{65}
在可以选择时，我们更喜欢用纳入标准，因为它直接描述研究中人的特征。一个好的通则是：排除标准越少越好。

许多临床试验通过纳入及排除标准，找出最可能从治疗中获益、最不可能发生不良反应的人。这可能意味着排除因既往治疗而更难治愈的患者；也有试验只招募对其他药物反应良好、看来也可能受益于新治疗的人。两种情况下，更好的选择都可能是纳入既往其他治疗失败、最需要新疗法的患者，因为他们很可能就是治疗的目标人群。

有些研究主要旨在认识疾病的生物学，例如特征和原因；另一些主要用于指导实际治疗决策，例如选择药物还是手术（5）。不同目的会影响选择标准。例如，主要关心睾酮补充精确生理作用的研究者，可能排除服用任何其他药物的女性，限定一个较窄的样本；另一位则可能只设少量排除条件，使结果有更广泛的适用性。

规定临床特征时，常常难以决定如何具体界定标准。例如，怎样操作“健康状况良好”这一条件？研究者也许决定不纳入任何自报疾病者，但这可能排除大量本来很适合研究问题的人，例如骨关节炎患者。因此，应说明每一项纳入与排除标准的理由，确认它们必要而恰当。

最后，应避免把可能由某个研究结局引起的因素作为选择标准；
\sourcepage{66}
例如，研究肺癌危险因素时，不应以近期体重下降作为纳入或排除条件。第 10 章“以共同结果为条件”一节会更详细讨论这一点。

\subsection{临床人群与有代表性的人群}
招募某种疾病的患者时，住院或门诊患者可能容易找到，但决定谁来医院或门诊的选择因素，可能产生重要影响。例如，学术医疗中心的专科门诊往往吸引病情严重者，使人对疾病特征和预后产生失真的印象。从社区医疗机构抽样可能更好。

样本常包括通过信件、电子邮件、互联网或大众媒体广告招募的人。但是，有些人比其他人更容易被联系到或更愿意参加，因而这样的样本不能充分代表一般人群，正如美国 2016 年和 2020 年大选前民调所显示的那样。真正\term{以人群为基础的样本}（population-based samples），例如能代表美国居民的美国国家健康与营养检查调查（National Health and Nutrition Examination Survey，NHANES），招募既困难又昂贵，却有助于指导公共卫生及临床实践（6）。

与其他地区的同行合作、使用 NHANES 或 Medicare 记录等已有数据集，或获得公共卫生部门、医疗机构及医疗保险公司的数据，可以扩大样本规模和多样性。这类数据已广泛用于临床研究，通常更能代表全国人群，而且“招募”所需时间更少（见第 16 章）。

由于不便、缺少社区参与乃至系统性种族主义等多种原因，少数族裔或贫困者等人群一直在研究中代表不足，使人难以足够精确地估计治疗或暴露对这些群体的作用。
\sourcepage{67}
这些不足可能并且已经造成健康差异；负责任的研究者应计划投入额外资源招募他们。

代表多样人群的样本，可使研究者考察不同群体间结果是否不同，例如男性与女性、不同种族或族裔之间。不过，除非这是具体研究目标，单项研究通常没有足够的亚组人数检出这些差异，除非差异很大。但是，按重要亚组报告结果，再与其他以同样方式报告结果的研究合并，就可能在系统综述或荟萃分析中识别差异。

\section{抽样}
符合选择标准的潜在人群往往太大，无法全部研究，因此必须从可接触人群选择一个子集作为样本。当然，如果研究使用现成的电子数据，就未必如此；取得必要批准后，研究者可能研究整个可接触人群。

抽样方法有两大类：一类明确规定总体中每个人被选中的概率，另一类则没有。因此分别称为\term{概率样本}（probability samples）及\term{非概率样本}（nonprobability samples）。从样本作出推断是否有效，取决于样本代表可接触人群的程度；对非概率样本，需要对此作出主观判断。

\subsection{概率样本}
概率抽样是统计推断的金标准，利用随机过程使可接触人群每个成员都有指定的入样机会（图 3.4）。这样能够有根据地估计样本中的现象在多大程度上代表可接触人群，并计算统计学显著性及置信区间。通常需要\term{抽样框}（sampling frame），即列出可接触人群中所有可能入组的人，
\sourcepage{68}
或下文将介绍的群组。概率抽样有几种形式。
\par
\begin{figure}[htbp]\centering\resizebox{\linewidth}{!}{\chthreesamples}
\caption*{\textbf{图 3.4}\quad 简单随机样本、分层随机样本和整群样本的例子。}
\end{figure}\par
\begin{description}[leftmargin=2em,labelindent=0pt,style=nextline]
\item[简单随机样本（simple random sample）。]通过随机过程从抽样框选择，每个人被选中的概率相同。例如，一名研究白内障手术结局的研究者，找出研究期间所在医疗系统接受该手术的所有患者。假定共有 2500 人，而她只有资源访谈 100 人。先给患者编上 1～2500 的号码，再为每人分配一个 0～1 的随机数；电子表格的 \texttt{RAND()} 函数即可完成。随后按随机数从小到大排序，取前 100 人。如果有人不参加，则用随机数接下来最小的人替代。
\item[系统样本（systematic sample）。]与简单随机样本不同，通过周期性程序选取，例如弗雷明汉研究在按地址排序的家庭名单中，每三个家庭取前两个。这种抽样容易受到总体自然周期性的影响，
\sourcepage{69}
也使研究者能够预测、甚至操纵谁进入样本。它在实施方面并不优于简单随机抽样，在临床研究中很少是更好的选择。
\item[分层随机样本（stratified random sample）。]先按性别、年龄等特征把总体分为亚组，再分别从各“层”随机抽样。可对分层子样本给予不同权重，过量抽取总体中较少见、但研究者特别关注的亚组。例如，研究妊娠期子痫前期发病率时，可按自报种族或族裔分层，每层抽取相同人数。较少见种族因此得到相对过量的代表，各层发病率估计便有相近的精度。
\item[整群样本（cluster sample）。]随机抽取总体中个人自然组成的群组。人群分布广泛，或无法实际列出并抽取所有成员时，这种方法很有用。例如，希望访谈全州登记系统中随机选出的肺癌患者时，可以先随机抽取诊断这些患者的医院，再访谈所选医院的病例，从而降低费用。
\end{description}
社区调查常用两阶段整群抽样：先从地图上列出的街区随机抽样，再通过第二次随机过程，从入选街区的地址中抽取子样本。整群抽样的缺点是，自然形成的群组通常比总体更同质，例如同一街区居民的社会经济地位常常相近。因此，校正群内相似性后的有效样本量略小于实际人数；数据分析以及样本量估计（第 6 章）都必须考虑群组结构。

\subsection{非概率样本}
\sourcepage{70}
概率样本虽然理想，但可能不便、昂贵，甚至因为无法建立抽样框而不可能取得。相比之下，不采用随机选择程序的非概率样本，通常更容易获得。
\begin{description}[leftmargin=2em,labelindent=0pt,style=nextline]
\item[连续样本（consecutive sample）。]按某种顺序纳入符合入组条件者，例如首先识别出的 50 人。优点是不必预先列出可接触人群所有成员，但在可推广性方面可能有缺点。招募期足以涵盖对研究问题重要的季节变化或其他时间变化时，这种方法效果最好。
\item[便利样本（convenience sample）。]因为研究者容易识别和招募而被选入的人。涉及与人接触的研究，都不可避免地考虑获取机会及可行性；但研究者应认识到，为方便而牺牲研究方法的周密程度，会降低科学价值。“便利样本”有时带有贬义。只要可行，概率抽样都更好。
\end{description}

\subsection{关于抽样的几点说明}
根据样本观察，利用描述统计及显著性检验推断可接触人群，前提是使用了概率抽样。但是，临床研究未必能够招募可接触人群的概率样本。因此，是否接受拟议抽样设计，通常需要研究者判断：对手头的研究问题，根据样本观察所得的结论，是否与研究真正概率样本所得结论相似？再往外一步，这些结论是否适用于目标人群？

不要把随机抽样与\term{随机化}或\term{随机分配}（randomization or random assignment）混淆。前者决定可接触人群哪些成员入组；后者则例如决定临床试验参与者接受有效治疗还是安慰剂（见第 11 章）。
\sourcepage{71}
两者都能在推断时获得使用随机数的益处，但很少有研究同时采用。一个显著的例外虽然不属临床研究，却是社交媒体公司：它们随机选择用户，并向用户展示随机分配的广告版本，以评估各种“吸引手段”的效果（7）。

许多期刊要求以图示说明抽样过程，包括招募及随访各阶段流失的潜在参与者人数。制订计划时就考虑这一要求，很有帮助。

\section{招募}
招募的可行性会影响抽样方法。招募有两个主要目标：其一，纳入能够充分代表可接触人群的样本，尽量避免因系统误差（偏倚）答错研究问题；其二，获得足够的样本量，尽量避免因随机误差（偶然性）得到错误答案。

\subsection{获得有代表性的样本}
招募代表性样本，始于设计阶段对目标人群、可接触人群及抽样类型的合理选择，终于实施阶段：防止在潜在参与者中错误应用入组标准，并随着研究进展加强有效的留存措施。

一个常见错误是，不同人群入组率不同，导致实际招募样本不能反映可接触人群。例如，自愿参加研究者通常更健康、受教育程度更高，也更少吸烟。黑人入组率往往低于其他人群，部分原因可能是过去某些公然带有种族主义的研究（见第 7 章）。用数字化方式接触参与者时，\term{自愿者偏倚}（volunteer bias）甚至可能更严重（8）。
\sourcepage{72}
要招募多样化样本，通常需要更多人与人的接触，着重建立研究人员与潜在参与者之间的关系。第 15 章社区参与式研究将进一步讨论招募多样且有代表性样本的综合方法。

\term{无应答}（nonresponse）也是一大问题，尤其是以选举民调为例的描述性研究。\footnote{本章讨论招募过程中无应答，主要关注以估计变量分布为主要目标的描述性研究。对随时间随访参与者的前瞻性研究，随访期间无应答是一项重大问题，尤其是干预本身可能改变应答率的临床试验（第 11 章）。}被选中者中同意入组的比例，即\term{应答率}（response rate），影响将入组样本视为可接触人群代表这一推断的有效性。难以联系到的人，以及联系后拒绝的人，往往与入组者不同。多高的无应答程度会损害可推广性，取决于问题性质及未应答原因。不过，即使无应答率仅 25\%——在许多场景已算不错——如果疾病本身影响应答可能性，也会严重扭曲患病率估计。

有时可以对部分无应答者补充收集资料，估计\term{无应答偏倚}（nonresponse bias）可能怎样影响结论，但这需要额外资源及 IRB 的特别批准。例如，通过社交媒体招募的研究，可以获取无应答者的公开资料；按门诊患者名单邮寄招募的研究，可随机抽取部分无应答者，提供金钱激励，请他们通过电话回答几个补充问题。

不过，最佳办法还是采用多次联系和多种方式（信件、电子邮件、电话），尽量减少无应答者。对已联系到的人，可提高研究的效率与吸引力，避免侵入性或令人不适的检查，通过宣传册及个别交谈减轻焦虑和不适，提供交通费用报销、告知检查结果等激励，并使用双语人员和译制问卷克服语言障碍，来减少拒绝参加。
\sourcepage{73}
另一种选择是为完成研究提供激励。

\subsection{招募足够多的参与者}
招募不足是临床研究最常见的问题之一。正如一些同行写道：
\begin{studyquote}
然而，随着时间推移，起初的招募热情逐渐消退。面对招募缓慢等研究的残酷现实，负面情绪常常占据上风。研究者可能感到绝望、自责、内疚、失败、无价值、孤独、挫折及亚临床抑郁；有时甚至出现偏执（9）。
\end{studyquote}
措辞虽然强烈，却反映了一个事实：规划研究时，最好假定符合入组标准且同意参加的人数将低于起初预计，有时甚至少上几倍。解决办法包括：通过预试验了解招募困难的实际程度；计划的样本量高于自己认为必要的数量；预先安排在需要时增加参与者的办法。招募进行中，应监测目标完成情况，并列出不足的原因。理解潜在参与者为何在各阶段流失，可以帮助制订减少流失的策略。

有时，招募对象是团队已经熟悉的人，例如在研究者门诊就诊、拟接受新疗法研究的患者。这时应公平说明参与机会，清楚呈现利弊。研究者必须认识到，自己的医师职责可能与研究利益冲突，由此产生伦理困境（第 7 章）。事实上，如果研究存在任何潜在风险或不便，而对象是自己的患者，几乎总应请一位无利害关系的同事取得其知情同意。
\sourcepage{74}
更常见的是，招募团队并不认识的潜在参与者。团队至少一人熟悉各种接触方式，会很有帮助。这些方式包括在工作场所或购物中心等公共地点筛查；利用驾驶执照持有人等名单大量寄发电子或纸质招募信；通过互联网或社交媒体广告；请临床医师转介；回顾病历；以及检查门诊及住院患者名单。部分方法涉及隐私侵犯，须经 IRB 考虑。

衰弱者或残障人士等特殊人群，可能需要家访或免费交通等额外招募方式。不能到门诊接受检查本身可能影响可推广性时，这类主动服务尤其重要。

事先取得有关组织支持，也可能有所帮助。例如，与医院管理者讨论门诊样本；与当地医学会商议支持向医师寄信；与社区领袖及县卫生部门计划社区筛查。可将书面支持函附在资助申请后。通过公众讲座、社区倡议活动，以及社交和大众媒体、传单、网站、邮寄宣传，也可营造有利的社区氛围。

不过，对于初学者，我们关于招募的最佳建议很简单：能够避免时，就避免。明确问题和计划后，首先查明其他团队是否已招募了足以回答问题的样本并完成相关测量，然后联系对方，看看是否可以获取和分析数据（见第 16 章），甚至在其后续研究访视中添加一项新测量。
\sourcepage{75}
这需要熟悉文献，并鼓起勇气联系可能尚不认识的研究者，却能省去数年努力。

\section{小结}
\begin{enumerate}
\item 大多数临床研究在理念及实践上，都是用样本代表总体。
\item 抽样的优点是效率：以较少时间和精力研究一个子集，就可推断较大的总体。缺点是引入误差：代表性不足，结果可能难以推广到目标人群；样本量不足，则难以充分减少偶然性的影响。
\item 设计样本时，先用具体纳入标准构思目标人群，明确适合研究问题的临床和人口学特征，再谨慎设置少量排除标准，排除在伦理上不宜或其他原因不适合研究的人。
\item 随后选择地域与时间上方便的可接触人群。
\item 下一步是设计从可接触人群抽样的方法。有完整名单时，概率抽样最好。纯粹因为方便而选人，可能得到有偏倚的样本。
\item 最后，应设计并实施招募及留存策略，使参与者能充分代表目标人群。
\end{enumerate}
\begin{references}
\item Framingham Heart Study. Epidemiological background and design: The Framingham Heart Study. \url{https://framinghamheartstudy.org/fhs-about/history/epidemiological-background}
\sourcepage{76}
\item Cummings SR, Nevitt MC, Browner WS, et al. Risk factors for hip fracture in white women. Study of Osteoporotic Fractures Research Group. N Engl J Med. 1995;332(12):767-773.
\item Cauley JA, Lui LY, Ensrud KE, et al. Bone mineral density and the risk of incident nonspinal fractures in black and white women. JAMA. 2005;293(17):2102-2108.
\item Orwoll E, Blank JB, Barrett-Connor E, et al. Design and baseline characteristics of the osteoporotic fractures in men (MrOS) study---a large observational study of the determinants of fracture in older men. Contemp Clin Trials. 2005;26(5):569-585.
\item Thorpe KE, Zwarenstein M, Oxman AD, et al. A pragmatic-explanatory continuum indicator summary (PRECIS): a tool to help trial designers. J Clin Epidemiol. 2009;62(5):464-475.
\item Centers for Disease Control and Prevention, NCHS. National Health and Nutrition Examination Survey. \url{https://www.cdc.gov/nchs/nhanes/index.htm}
\item Bakshy E, Eckles D, Yan E, et al. Social influence in social advertising: evidence from field experiments. In: Proceedings of the 13th ACM Conference on Electronic Commerce (EC '12). ACM; 2012:146-161. \url{https://research.fb.com/wp-content/uploads/2016/11/social-influence-in-social-advertising-evidence-from-fieldexperiments.pdf}
\item Guo X, Vittinghoff E, Olgin JE, Marcus GM, Pletcher MJ. Volunteer participation in the Health eHeart Study: a comparison with the US population. Sci Rep. 2017;7(1):1956.
\item Patel M, Doku V, Tennakoon L. Challenges in recruitment of research participants. Adv Psych Treatment. 2003;9:229-238.
\end{references}

\sourcepage{77}
\chapterappendix{3A}{第 3 章习题：选择研究参与者}
\begin{enumerate}
\item 研究者关注的问题是：“哪些因素使人开始吸烟？”她决定采用高中生横断面样本，邀请自己所在郊区高中的十一年级学生参加，并研究其中的自愿者。
\begin{enumerate}
\item 讨论这一样本对于所关心目标人群的适合程度。
\item 假定她为避免自愿者偏倚，计划从整个十一年级随机抽取 25\%，但实际样本中女生占 70\%。已知全年级男女人数相等，那么性别分布不成比例，就是抽样中的误差。它可能由随机误差、系统误差，还是两者造成？解释你的答案。
\end{enumerate}
\item 研究者准备调查摇滚音乐会观众，了解他们对演出期间佩戴耳塞保护听力的态度。请指出以下选择填写简短问卷者的抽样方法名称，并评论其可行性，以及结果能否推广到所有摇滚音乐会观众。
\begin{enumerate}
\item 每名观众入场时，在研究者手机上掷一个虚拟骰子；掷出 6 的人获邀填问卷。
\item 每名观众入场时掷虚拟骰子；男性掷出 1、女性掷出偶数者获邀。
\item 门票在售票处按号码顺序出售，票号末位为 1 的观众获邀。
\item 观众全部就座后，从每排对应一张、已洗匀的卡片中随机抽出五排，邀请这五排所有观众。
\item 邀请最先入场的 100 名观众。
\item 部分票邮购，部分在演出前由售票处出售。只要排队购票者达到五人或以上，就邀请队尾最有时间的人。
\sourcepage{78}
\item 演出后观众开始离场时，邀请那些看来愿意而且能够填问卷的人。
\end{enumerate}
\item Edwards 等（1）报告了 5 岁以下儿童人偏肺病毒（human metapneumovirus，HMPV）的感染负担。对象来自辛辛那提、纳什维尔及纽约州罗切斯特周边各县，是 2003～2009 年每年 11 月至次年 5 月因急性呼吸道疾病或发热就医的儿童。研究于周日至周四招募同意的住院患儿，每周一两天招募门诊患儿，每周一至四天招募急诊患儿。作者把各研究地点检测阳性的比例，与全国门诊医疗调查及全国医院门诊医疗调查中因急性呼吸道疾病或发热就诊的人群频率结合，估计美国 HMPV 的总体负担。他们估计，每年每 1000 名儿童中，HMPV 导致 55 次门诊就诊和 13 次急诊就诊。
\begin{enumerate}
\item 目标人群是什么？
\item 可接触人群是什么？由其推广到目标人群有多合适？
\item 采用了什么抽样方法？由其推广到可接触人群有多合适？
\item 概括说明，计算所估计 HMPV 发生率的置信区间时，应如何考虑这种抽样方法。
\end{enumerate}
\end{enumerate}
\begin{references}
\item Edwards KM, Zhu Y, Griffin MR, et al. Burden of human metapneumovirus infection in young children. N Engl J Med. 2013;368:633-643.
\end{references}
